\chapter{All-Genus Arithmetic Koszul Duality}
\label{v4-arith:all-genus-koszul}

Chapter~\ref{v4-arith:rh-reformulation} establishes the \(g=0\) case: the
internal structural identity
\(\kappa^{\mathrm{Ar}}_{\mathsf{G}}+(\kappa^{!})^{\mathrm{Ar}}_{\mathsf{G}}=0\)
of \(\mathcal{V}^{\mathrm{prim}}\) (the Hochschild trace class of the identity
on the global arithmetic Iwahori--Hecke category of \(GL_2\)) is equivalent
to the Riemann functional equation \(\xi(s)=\xi(1-s)\), via adelic Poisson
summation and rank-\(2\) Satake parameters
(\Cref{v4-arith:thm:P1-resolution}). The question this chapter answers:
for an arithmetic datum of rank \(2g\) with \(g\geq 0\), which internal
identity is equivalent to the corresponding motivic functional equation?

The answer is the same identity
\(\kappa^{\mathrm{Ar},g}_{\mathsf{G}}+(\kappa^{!})^{\mathrm{Ar},g}_{\mathsf{G}}=0\),
holding at every \(g\). The genus does not deform the Koszul sum; it
enters as the rank of the Heisenberg oscillators at each place, the
archimedean \(\Gamma\)-factor \(\Gamma_{\mathbb{C}}(s)^g\), and the
motivic weight fixing the functional-equation center.

\section{The Arithmetic Genus-g Datum}
\label{v4-arith:ag:datum}

\begin{definition}[arithmetic genus-\(g\) datum]
\label{v4-arith:ag:defn:datum}
\ClaimStatusDefinitional.
An \emph{arithmetic genus-\(g\) datum} \(\mathcal{D}_g\) is a triple
\((\mathcal{B}_g,\mathcal{E}_g,\varphi_g)\) specifying:
\begin{enumerate}
\item \(\mathcal{B}_g\), an arithmetic base: one of
\begin{enumerate}
\item Arakelov-compactified \(\overline{\mathrm{Spec}\,\mathcal{O}_K}\) for
\(K\) totally imaginary of degree \(2g\) (Window~II);
\item a proper flat regular arithmetic surface \(X\to\mathrm{Spec}\,\mathcal{O}_K\)
whose generic fibre \(X_K\) is a smooth projective curve of genus \(g\),
Arakelov-compactified via the archimedean Kähler form
\(\mu_{\mathrm{Ar},g}\) of Faltings 1984 (Windows~I, V);
\item a Shimura variety \(\mathrm{Sh}_K(G,X)\) of a Shimura datum whose
complex fibre is a moduli of abelian \(g\)-folds (Window~VI).
\end{enumerate}
\item \(\mathcal{E}_g\), a rank-\(2g\) arithmetic cohomology object:
\(H^{1}_{\mathrm{prim}}(\mathcal{B}_g)\) carrying Galois action (étale or
crystalline), with Hodge structure of weight \(1\) at each archimedean place:
\(\mathcal{E}_g\otimes\mathbb{C}=\mathcal{E}_g^{1,0}\oplus\mathcal{E}_g^{0,1}\)
with \(\dim\mathcal{E}_g^{1,0}=\dim\mathcal{E}_g^{0,1}=g\) per place.
\item \(\varphi_g\), a Frobenius-compatible endomorphism on \(\mathcal{E}_g\)
with square equal to the appropriate motivic involution: \(\varphi_g\) is
either
\begin{enumerate}
\item the archimedean Frobenius \(\Phi_\infty\) of Deninger 1998 acting on
the primitive Minkowski cohomology (Window~II);
\item the prismatic Frobenius \(\varphi\) of Bhatt--Scholze 2022 on the
universal rank-\(2g\) \(F\)-crystal over the ordinary crystalline locus of the
Emerton--Gee stack \(\mathcal{X}_{2g}\) (Window~III);
\item the Koszul duality operator \(J^{\mathrm{Koz}}\) on chiral cohomology
satisfying \((J^{\mathrm{Koz}})^{2}=\tau\) for \(\tau\) the motivic
functional-equation involution (Window~V).
\end{enumerate}
\end{enumerate}
At \(g=0\) the datum degenerates to the base \(\overline{\mathrm{Spec}\,\mathbb{Z}}\)
with \(\mathcal{E}_0\) void and \(\varphi_0=\mathrm{id}\); the functional
equation reduces to \(\xi(s)=\xi(1-s)\) and \Cref{v4-arith:ag:thm:all-genus}
specialises to \Cref{v4-arith:thm:P1-resolution}.
\end{definition}

\begin{remark}[Frobenius as almost-complex substitute]
\label{v4-arith:ag:rem:almost-complex}
An endomorphism \(J:TX\to TX\) satisfying \(J^2=-\mathrm{id}\) cannot exist
on an arithmetic scheme: schemes over \(\mathrm{Spec}\,\mathbb{Z}\) carry no
real tangent bundle, and on a compact real \(2\)-manifold
Newlander--Nirenberg (1957) makes integrability of \(J\) automatic, so no
genuinely almost-complex geometry survives.  The Frobenius
\(\varphi_g\) of \Cref{v4-arith:ag:defn:datum} is the correct
arithmetic substitute: on archimedean primitive cohomology it acts as
Hodge rotation by \(\pi/2\); on the prismatic \(F\)-crystal ordinary
stratum it is the Newton--Hodge-polygon-aligned slope-decomposition
Frobenius (Katz 1981, \emph{Ast\'erisque}~63); on chiral cohomology it
is the Koszul duality operator \(J^{\mathrm{Koz}}\) whose square equals
the functional-equation involution \(\tau\).  No non-integrable
decoration is introduced.
\end{remark}

\section{Six Convergent Windows}
\label{v4-arith:ag:windows}

The datum \(\mathcal{D}_g\) admits six distinct geometric and automorphic
presentations, each extracting a different structural feature of the
rank-\(2g\) arithmetic cohomology and the corresponding motivic functional
equation.

\subsection{Window I: Arakelov fibration}
\label{v4-arith:ag:w1-arakelov}

\(\mathcal{B}_g^{\mathrm{I}}=X\to\mathrm{Spec}\,\mathcal{O}_K\), proper flat
regular with generic fibre of genus \(g\), Arakelov-compactified.
\(\mathcal{E}_g^{\mathrm{I}}=H^{1}_{\mathrm{B}}(X(\mathbb{C}),\mathbb{Z})\)
of rank \(2g\) per archimedean place. \(\varphi_g^{\mathrm{I}}\) the integrable
complex structure \(J_\infty\) at each archimedean place. Arakelov 1974,
Faltings 1984, Gillet--Soulé 1990 (Publ.~IHES) supply the arithmetic
intersection pairing on \(\widehat{\mathrm{CH}}^1(X)\); Faltings--Hriljac
1984 the arithmetic Hodge index. Genus \(g\) enters via \(\pi_*\omega_{X/S}\)
of rank \(g\), the \(\Gamma\)-factor \(\Gamma_{\mathbb{C}}(s)^g\) of the
Hasse--Weil \(L\)-function, the coefficient \(2g-2\) in \(\deg\omega_{X_K/K}\),
and the \(\delta\)-invariant correction in the arithmetic Noether formula
\(\widehat{\deg}\det R\pi_*\omega=\tfrac{1}{12}(\widehat{\omega}^2+\delta_{\mathrm{Ar}})\)
(Faltings 1984, Thm~5).

\subsection{Window II: Minkowski torus / Bost--Connes--Marcolli}
\label{v4-arith:ag:w2-minkowski}

\(\mathcal{B}_g^{\mathrm{II}}=\overline{\mathrm{Spec}\,\mathcal{O}_K}\) for
\(K\) totally imaginary of degree \(2g\), with Minkowski lattice
\(\Lambda_K=j(\mathcal{O}_K)\subset\mathbb{C}^g\) and Minkowski torus
\(T_K=\mathbb{C}^g/\Lambda_K\). \(\mathcal{E}_g^{\mathrm{II}}=
H^{1}(T_K,\mathbb{Z})\cong\Lambda_K^\vee\) of \(\mathbb{Z}\)-rank \(2g\);
the primitive part equals the full \(H^1\) here since \(T_K\) is an abelian
variety.
\(\varphi_g^{\mathrm{II}}=\Phi_\infty^{\mathrm{tw},K}\), the half-Tate-twisted
archimedean Frobenius acting as \(\sqrt{-1}\) on each complex embedding
pair; equivalently the \(\mathbb{C}^g\)-structure on \(\Lambda_K\otimes\mathbb{R}\).
The genus is \(g=r_2(K)=\tfrac{1}{2}[K:\mathbb{Q}]\).
Bost--Connes--Marcolli 2004 (\emph{Selecta Math.}~11) constructs the
\(C^*_K\) quantum statistical system whose KMS\(_\beta\) partition function
is \(\zeta_K(\beta)\); Ha--Paugam 2005 and Laca--Neshveyev 2011 classify
its KMS states.

\subsection{Window III: Ordinary crystalline Emerton--Gee stack}
\label{v4-arith:ag:w3-fargues-scholze}

\(\mathcal{B}_g^{\mathrm{III}}=\mathcal{X}_{2g}^{\mathrm{crys,ord}}\), the
ordinary crystalline locus of the Emerton--Gee stack \(\mathcal{X}_{2g}\)
(Emerton--Gee arXiv:1908.07185) as a \(2\)-dim Artin stack over
\(\mathrm{Spec}\,\mathbb{Z}_p\), with prismatic upgrade via Bhatt--Scholze 2022
(\emph{Ann.~Math.}~196). A point is a rank-\(2g\) ordinary crystalline
\(\varphi\)-module with Hodge--Tate weights \(\{0,1\}^g\); equivalently a
rank-\(2g\) \(F\)-crystal with Newton polygon \((0,1)^g\). \(\mathcal{E}_g^{\mathrm{III}}\)
is the universal rank-\(2g\) bundle; \(\varphi_g^{\mathrm{III}}\) is the
prismatic Frobenius splitting \(\mathcal{E}_g^{\mathrm{III}}\) as
\(\mathcal{O}(0)^g\oplus\mathcal{O}(1)^g\) on the ordinary stratum.
Archimedean tilting relates this to the Hodge decomposition
\(H^1(X(\mathbb{C}),\mathbb{C})=H^{1,0}\oplus H^{0,1}\).
Caraiani--Scholze 2017--2019 (arXiv:1511.02418; arXiv:1709.02590) genericity
pins the middle-degree Galois representation to rank \(2g\) on compact
unitary Shimura surfaces of signature \((1,2)\) with generic Hecke eigenclass.

\subsection{Window IV: Hodge--Arakelov / arithmetic Teichmuller}
\label{v4-arith:ag:w4-mochizuki}

\(\mathcal{B}_g^{\mathrm{IV}}=[C,\vartheta]\in\mathcal{T}_g^{\mathrm{arith}}(\mathcal{O}_K)\),
a closed point of the arithmetic Teichmüller space carrying a level-\(n\)
theta-structure \(\vartheta\) on the Jacobian \(J(C)\). \(\mathcal{E}_g^{\mathrm{IV}}\)
is the \(n^g\)-dim Jet module over the theta-structure locus;
\(\varphi_g^{\mathrm{IV}}\) is Mochizuki's Hodge--Arakelov comparison
isomorphism (Mochizuki 1999 \emph{Publ.~RIMS}; arXiv:math/0304054). The
comparison is theorem-grade for \(g=1\); conjectural for \(g\geq 2\). The
Steinberg relation \(\{p,1-p\}=0\in K_2^M(\mathbb{Q})\) is preserved under
every Frobenioid indeterminacy (anabelian invariant via
Matsumoto 1969, Milnor 1970). The na\"ive speculation conflates the crystalline--\'etale comparison
isomorphism with an almost-complex structure \(J\) on
\(\mathrm{Spec}\,\mathbb{Z}\); the correct arithmetic object is the
Hodge--Arakelov decoration \(\varphi_g^{\mathrm{IV}}\) on the dense
theta-generic open locus of \(\mathcal{T}_g^{\mathrm{arith}}\)
(\Cref{v4-arith:ag:rem:almost-complex}).

\subsection{Window V: \(E_2\)-chiral factorisation on the arithmetic surface}
\label{v4-arith:ag:w5-bar-cobar}

\(\mathcal{B}_g^{\mathrm{V}}=\mathcal{S}_g^{\mathrm{Ar}}\), the arithmetic
surface of Window~I viewed through its \(E_2\)-chiral factorisation
structure. \(\mathcal{E}_g^{\mathrm{V}}=\mathcal{V}^{\mathrm{prim,surf},g}:=
\mathrm{HH}_\bullet(\mathrm{Hk}^{\mathrm{Iw,surf}}_{GL_{2g},\mathrm{glob}})\),
the Hochschild trace class of the identity on the \(E_2\)-monoidal global
Iwahori--Hecke category of \(GL_{2g}\) on the arithmetic surface. The \(E_2\)
structure arises via \(E_1\otimes E_1=E_2\) Dunn additivity (Lurie
\emph{Higher Algebra}~5.1.2.2) from the horizontal \(E_1\) (1-dim base
factorisation on \(\overline{C}\)) and the vertical \(E_1\) (fibre
factorisation on \(X_K\)); the bar-cobar construction on this \(E_2\) datum
lifts to the 2-dim Parshin--Gersten complex with double residue
\((\partial_1,\partial_2)\) (Parshin 1976; Kato 1980; Beilinson 1980).
\(\varphi_g^{\mathrm{V}}=J^{\mathrm{Koz}}\), the Koszul duality operator
on chiral cohomology with \((J^{\mathrm{Koz}})^2=\tau\) (the
functional-equation involution). The genus enters via rank-\(2g\) Fock
oscillators per place, sourced from the \(\ell\)-adic Tate module
\(T_\ell(\mathrm{Jac}(X_K))\); the holomorphic central charge is \(c=g\),
with \(g\) additional antiholomorphic oscillators, giving total Fock
rank \(2g\).

\subsection{Window VI: Kudla--Rapoport / Shimura--Siegel Shimura variety}
\label{v4-arith:ag:w6-shimura}

\(\mathcal{B}_g^{\mathrm{VI}}=\mathrm{Sh}_K(G,X)\) a Shimura variety with
\(G\in\{\mathrm{GSp}_{2g},\,\mathrm{Res}_{F/\mathbb{Q}}\mathrm{GL}_2\text{ for
}[F:\mathbb{Q}]=g,\,\mathrm{GU}(g,1)\}\). \(\mathcal{E}_g^{\mathrm{VI}}\) the
middle-degree cohomology decomposed under the Hecke algebra into rank-\(2g\)
(or \(2g+1\), \(4g\)) Galois blocks per automorphic \(\pi\) of \(G(\mathbb{A})\).
\(\varphi_g^{\mathrm{VI}}\) the integrable complex structure sourced from the
Deligne torus \(h:\mathbb{S}\to G_{\mathbb{R}}\). The Kudla--Rapoport
generating series \(\widehat{\phi}(\tau)=\sum_m\widehat{\mathcal{Z}}(m)q^m\)
is a genus-\(g\) Siegel modular form valued in \(\widehat{\mathrm{CH}}^1(\mathrm{Sh})_{\mathbb{C}}\)
(Kudla 1997 \emph{Ann.~Math.}~146; Kudla--Rapoport 2011
\emph{Algebra Number Theory}~5; Li--Liu 2021 arXiv:2101.09485). The \(g\)-fold arithmetic
intersection of special cycles realises Fourier coefficients of a
genus-\(g\) Siegel Eisenstein series' central derivative. The Langlands-dual
group \({}^LG=\mathrm{GSpin}_{2g+1}\) carries the \((2g+1)\)-dim spin \(L\)-function.

\section{Falsified Naive Readings}
\label{v4-arith:ag:falsified}

Six naive readings of the speculative phrase ``genus-\(g\) arithmetic
surface with almost-complex structure'' each fail at a specific technical
joint; the failures are recorded in \Cref{v4-arith:app:B:genus-g-session}.

\begin{enumerate}
\item \emph{\(J^2=-\mathrm{id}\) on the tangent bundle of an arithmetic
scheme.} Schemes over \(\mathrm{Spec}\,\mathbb{Z}\) carry no real tangent
bundle; the phrase ``almost-complex'' is a category error at this level.
The correct arithmetic substitute is the Frobenius \(\varphi_g\) on
primitive cohomology (\Cref{v4-arith:ag:rem:almost-complex}).
\item \emph{Arithmetic Euler characteristic \(=1-g\).}
Gillet--Soul\'e 1992 arithmetic Riemann--Roch introduces
Bismut--Gillet--Soul\'e analytic torsion; the arithmetic analog of \(1-g\)
is the Faltings \(\delta\)-invariant via the arithmetic Noether formula,
not a simple integer.
\item \emph{K\"unneth factorisation of the arithmetic surface.}
Bost--K\"unnemann 2010 (\emph{Adv.~Math.}~223) shows the product formula
has non-vanishing Bott--Chern archimedean correction terms; \(X\) does
not factor as ``genus-\(g\) curve
\(\times_{\mathrm{Spec}\,\mathbb{Z}}\mathrm{Spec}\,\mathbb{Z}\)'' in
arithmetic Chow.
\item \emph{\(\overline{\mathrm{Spec}\,\mathbb{Z}}\) itself has finite genus.}
Deninger's arithmetic-dynamical genus of \(\mathrm{Spec}\,\mathbb{Z}\) is
\(+\infty\) (countably many critical zeros of \(\zeta\)). Genus enters the
arithmetic datum only through coefficient enhancement to a totally
imaginary number field of degree \(2g\).
\item \emph{Conflation of curve-genus \(g_c\) with Siegel-modular genus
\(g_s\).} For a modular curve \(X_0(N)\): \(g_c=g_0(N)\) and \(g_s=1\).
For the Siegel threefold \(\mathcal{A}_2\): \(g_c\) is undefined, \(g_s=2\).
The genus \(g\) of the datum \(\mathcal{D}_g\) is the symplectic rank of the
Galois representation divided by \(2\); in the Shimura setting this is
\(g_s\), not \(g_c\).
\item \emph{\(\kappa+\kappa^!=2-2g\) as Euler-characteristic shift.}
The Euler characteristic \(2-2g\) is a topological invariant of the
Riemann surface \(X(\mathbb{C})\), not a Koszul invariant of the chiral
datum. The identity \(\kappa^{\mathrm{Ar},g}+(\kappa^{!})^{\mathrm{Ar},g}=0\)
holds at every \(g\); the genus enters through the Fock rank, the
archimedean \(\Gamma\)-factor, and the motivic weight of the
functional-equation center.
\end{enumerate}

\section{All-Genus Arithmetic Koszul Duality}
\label{v4-arith:ag:load-bearing}

Let \(\mathcal{D}_g=(\mathcal{B}_g,\mathcal{E}_g,\varphi_g)\) be an arithmetic
genus-\(g\) datum (\Cref{v4-arith:ag:defn:datum}). Attach to \(\mathcal{D}_g\)
the arithmetic Heisenberg \(\mathcal{V}^{\mathrm{prim},g}_{\mathrm{Heis}}\)
via the rank-\(2g\) Fock on \(\mathcal{E}_g\) at each place, and the Koszul
dual via the genus-\(g\) Bost--Connes--Marcolli system of Consani--Marcolli
2006 (\emph{J.~Reine Angew.~Math.}~597).

\begin{theorem}[all-genus arithmetic Koszul duality]
\label{v4-arith:ag:thm:all-genus}
\ClaimStatusProvedHereConditional \emph{(on the residual obligations of
\Cref{v4-arith:ag:residual})}. For every \(g\geq 0\) and every arithmetic
genus-\(g\) datum \(\mathcal{D}_g\), the arithmetic Heisenberg Koszul
self-duality identity
\begin{equation}
\label{v4-arith:ag:eq:master}
\kappa^{\mathrm{Ar},g}_{\mathsf{G}}(\mathcal{D}_g)
+(\kappa^{!})^{\mathrm{Ar},g}_{\mathsf{G}}(\mathcal{D}_g)=0
\end{equation}
is equivalent to the motivic functional equation
\begin{equation}
\label{v4-arith:ag:eq:fe}
\Lambda_{\mathcal{D}_g}(s)
=\epsilon_{\mathcal{D}_g}\cdot\Lambda_{\mathcal{D}_g}\bigl(w_{\mathcal{D}_g}-s\bigr),
\end{equation}
where \(\Lambda_{\mathcal{D}_g}(s)\) is the completed \(L\)-function attached to
\(\mathcal{D}_g\) and \(w_{\mathcal{D}_g}\in\mathbb{Z}_{\geq 1}\) is the
motivic weight center, depending on the window but not on \(g\) alone.
Forward implication \eqref{v4-arith:ag:eq:master}\(\Rightarrow\)\eqref{v4-arith:ag:eq:fe}
holds unconditionally via Euler-product identification, Consani--Marcolli
Koszul dual, and Weil--Bloch--Kato adelic Poisson. Full biconditional is
conditional on the 2-dim arithmetic Positselski equivalence at
\(\Lambda_2=\mathbb{Z}_p[[\Gamma_1\times\Gamma_2]]\)-coefficients (Window~V)
plus Koszul--Petersson compatibility at rank \(g\) (extending
Chapter~\ref{v4-arith:kp-compat}).
\end{theorem}

The four rows of \eqref{v4-arith:ag:eq:fe}, indexed by window, are
recorded in \Cref{v4-arith:ag:table:rows}.

\begin{table}[h]
\caption{Rows of the all-genus arithmetic Koszul duality.}
\label{v4-arith:ag:table:rows}
\centering
\begin{tabular}{llll}
\toprule
Window & \(\Lambda_{\mathcal{D}_g}(s)\) & \(w_{\mathcal{D}_g}\) & Motivic weight \\
\midrule
II & \(\xi_K(s)=|d_K|^{s/2}\Gamma_{\mathbb{C}}(s)^g\zeta_K(s)\) & \(1\) & \(0\) (Dedekind) \\
I, V & \(\Lambda(s,X_K)=N_{X_K}^{s/2}\Gamma_{\mathbb{C}}(s)^g L(s,X_K)\) & \(2\) & \(1\) (Hasse--Weil) \\
VI & \(\Lambda(s,\pi,\mathrm{spin})\) for \(\pi\) on \(\mathrm{GSp}_{2g}\) & \(2g+1\) & \(g\) (Siegel spin) \\ % RECTIFICATION-FLAG: Gate 1: w=2g+1 is the motivic weight of the spin representation of GSpin_{2g+1}; verify against Arthur 2013 Endoscopic Classification and specific spinor L-function normalisation before final inscripton.
III, IV & spectral refinements & indirect & refines above \\
\bottomrule
\end{tabular}
\end{table}

\begin{remark}[weight structure]
\label{v4-arith:ag:rem:weight}
The functional-equation center \(w_{\mathcal{D}_g}\) is set by the motivic
weight of the row, not by \(g\) itself. Window~II (Dedekind \(\zeta_K\) with
archimedean \(\Gamma_{\mathbb{C}}^g\) factor) is weight \(0\); the center
stays at \(s\mapsto 1-s\) regardless of \(g\), since \(\zeta_K\) for \(K\) of
degree \(2g\) has the same \(1-s\) symmetry as \(\zeta\) for any number field.
Windows~I, V (arithmetic surface with generic fibre of genus \(g\)) are
weight \(1\) on the Jacobian's \(H^1\)-motive; the center is \(s\mapsto 2-s\).
Window~VI (Siegel spin of genus-\(g\) cusp form) is weight \(g\); the center
\(s\mapsto 2g+1-s\) scales with \(g\) linearly (Arthur 2013,
\emph{Endoscopic Classification}). All three rows share
\(\kappa+\kappa^!=0\); the genus grading lives in the \(\Gamma\)-factor, the
Fock rank, and the weight of the row.
\end{remark}

\begin{proof}[Sketch]
Four steps, parallel to \Cref{v4-arith:thm:P1-resolution} at \(g=0\).

\emph{Step 1: Euler-product identification.} At each prime \(p\) of good
reduction for \(\mathcal{D}_g\), writing \(\alpha_1^{(p)},\ldots,\alpha_{2g}^{(p)}\)
for the Frobenius eigenvalues on \(\mathcal{E}_{g,p}\), the rank-\(2g\)
Heisenberg Fock partition function at inverse temperature \(\beta=s\) factors as
\(\prod_{j=1}^{2g}(1-\alpha_j^{(p)}p^{-s})^{-1}
=\det(1-p^{-s}\varphi_g|\mathcal{E}_{g,p})^{-1}\).
The global Euler product (over primes of good reduction, with local factors
at bad primes defined as in Serre 1970) gives \(\Lambda_{\mathcal{D}_g}(s)\)
in the respective window.

\emph{Step 2: Consani--Marcolli genus-\(g\) Koszul dual.}
Consani--Marcolli 2006 (\emph{J.~Reine Angew.~Math.}~597) constructs the
rank-\(2g\) Bost--Connes--Marcolli system with KMS partition function
\(\Lambda_{\mathcal{D}_g}(\beta)\) for inverse temperature
\(\beta>w_{\mathcal{D}_g}/2\) and meromorphic continuation. The Koszul dual
of \(\mathcal{V}^{\mathrm{prim},g}_{\mathrm{Heis}}\) is this genus-\(g\) system.

\emph{Step 3: rank-\(g\) Positselski pairing.} Arithmetic Positselski
(Positselski 2011 Mem.~AMS~996; Chapter~\ref{v4-arith:arith-positselski})
at rank \(g\), extended to \(\Lambda_2\)-coefficients for Window~V. Pairs the
bar coalgebra of the rank-\(2g\) Heisenberg with the cobar algebra of the
genus-\(g\) Bost--Connes system; total weight vanishes modulo the 2-dim
\(\mu\)-invariant vanishing condition (extending Ferrero--Washington 1979
\emph{Ann.~Math.}~109 for abelian at rank \(1\)).

\emph{Step 4: adelic Fourier symmetry \(s\mapsto w_{\mathcal{D}_g}-s\).}
Tate 1950 at \(g=0\) (Window~II with \(K=\mathbb{Q}\)); Weil 1967 \emph{Basic
Number Theory} for totally imaginary base (Window~II general); Bloch--Kato
1990 and Fontaine--Perrin-Riou 1994 for weight-\(1\) motives on Jacobians
(Windows~I, V); Arthur 2013 for Siegel spin \(L\)-functions (Window~VI).
Each window's adelic Fourier symmetry, applied to the adelic theta
distribution on \(\mathcal{E}_g\), gives \eqref{v4-arith:ag:eq:fe} with the
corresponding \(w_{\mathcal{D}_g}\).

Steps~1, 2, 4 close unconditionally. Step~3 is the residual obligation.
Forward implication is therefore unconditional; full biconditional
requires rank-\(g\) Positselski.
\end{proof}

\section{Reduction to \texorpdfstring{\(g=0\)}{g=0} and the F-RH Axis}
\label{v4-arith:ag:reduction}

At \(g=0\):
\begin{itemize}
\item Window~II: \(K=\mathbb{Q}\), \(\xi_K(s)=\xi(s)\), \(\Gamma_{\mathbb{C}}(s)^0=1\),
\(\zeta_K=\zeta\). Functional equation \(\xi(s)=\xi(1-s)\).
\item Windows~I, V: \(X_K=\mathbb{P}^1/K\), Jacobian trivial, \(\Lambda(s,X_K)\)
degenerates to \(\xi(s)\) with center \(s\mapsto 1-s\); the weight-\(1\) motive
is void and the weight-\(0\) Dedekind row takes over.
\item Window~VI: \(\mathcal{A}_0=\mathrm{pt}\), Siegel spin \(L\)-function
collapses to \(\zeta\), center \(s\mapsto 1-s\).
\end{itemize}

All four rows collapse to \(\xi(s)=\xi(1-s)\), recovering
\Cref{v4-arith:thm:P1-resolution} as the \(g=0\) specialisation of
\Cref{v4-arith:ag:thm:all-genus}.

\begin{corollary}[all-genus F-RH]
\label{v4-arith:ag:cor:all-genus-FRH}
\ClaimStatusConditional \emph{(on \Cref{v4-arith:ag:thm:all-genus} plus
rank-\(g\) extensions of the P2, P3, VB pillars of
Chapter~\ref{v4-arith:rh-reformulation})}. For every \(g\geq 0\) and every
arithmetic genus-\(g\) datum \(\mathcal{D}_g\) with rank-\(g\) Koszul--Petersson
compatibility and rank-\(g\) arithmetic Zhu modularity, the non-trivial
zeros of \(\Lambda_{\mathcal{D}_g}(s)\) lie on the self-dual locus
\(\mathrm{Re}(s)=w_{\mathcal{D}_g}/2\).
\end{corollary}

\begin{remark}[reach]
\label{v4-arith:ag:rem:reach}
\Cref{v4-arith:ag:cor:all-genus-FRH} subsumes: the Riemann Hypothesis
(\(g=0\), Window~II); the Hasse--Weil Riemann Hypothesis for Jacobians of
curves of arbitrary genus (Windows~I, V); the generalised Riemann
Hypothesis for Dedekind zeta of totally imaginary fields (Window~II at
higher \(g\)); the Siegel-spin Riemann Hypothesis for automorphic
\(L\)-functions of cusp forms on \(\mathrm{GSp}_{2g}(\mathbb{A}_{\mathbb{Q}})\)
(Window~VI). The minimal sufficient path is a rank-\(g\) upgrade of
Chapter~\ref{v4-arith:rh-reformulation}'s three-pillar implication
structure plus \Cref{v4-arith:ag:thm:all-genus}.
\end{remark}

\section{Residual Obligations}
\label{v4-arith:ag:residual}

\Cref{v4-arith:ag:thm:all-genus} reduces to ten residual obligations,
each extending a Chapter-level \(g=0\) result to arbitrary \(g\).

\begin{enumerate}
\item \emph{2-dim arithmetic Positselski equivalence} with
\(\mu_1=\mu_2=0\) on \(\Lambda_2=\mathbb{Z}_p[[\Gamma_1\times\Gamma_2]]\),
extending Chapter~\ref{v4-arith:arith-positselski}. Scope: 60--100
pages. Precedents: Nekov\'a\v{r} 2006 (Astérisque~310); Coates--Sujatha
2006 \emph{Cyclotomic Fields and Zeta Values}.
\item \emph{Rank-\(g\) Koszul--Petersson compatibility} replacing the
scalar Petersson inner product with the genus-\(g\) Petersson pairing on
weight-\((k,g)\) automorphic forms, extending Chapter~\ref{v4-arith:kp-compat}.
Scope: 15--25 pages.
\item \emph{2-dim Parshin--Gersten bar construction} on the arithmetic
surface with double residue \((\partial_1,\partial_2)\), extending
Chapter~\ref{v4-arith:arakelov-bar}. Scope: 40--60 pages. Precedents:
Parshin 1976, Kato 1980, Beilinson 1980.
\item \emph{Rank-\(g\) arithmetic Ben-Zvi--Nadler} with target stack
\(\Sigma_g^{\mathrm{Ar}}\) replacing \(B\widehat{\mathbb{Z}}\), extending
Chapter~\ref{v4-arith:arith-bzn}. Scope: 80--120 pages.
\item \emph{\(GL_{2g}\)-local geometric Langlands.} Fargues--Scholze
arXiv:2102.13459 and Emerton--Gee arXiv:1908.07185 hold for arbitrary
reductive \(G\); the categorical comparison conjecture at \(GL_{2g}\) is
the rank-\(g\) scaling of \Cref{v4-arith:conj:F1-global-AG}. Scope:
determined by F1 scaled to rank \(2g\).
\item \emph{Perfectoid \(E_2\)-enhancement at rank \(g\)} via BD
Grassmannians for \(GL_{2g}/\mathbb{Q}_p\), extending F2
(\Cref{v4-arith:conj:F2-perfectoid-E2}). Scope: 60--80 pages.
\item \emph{Rank-\(g\) arithmetic Virasoro bootstrap} with central
charge \(c^{\mathrm{Ar},g}=g\) on the primary subspace of the rank-\(g\)
archimedean Fock, extending Chapter~\ref{v4-arith:virasoro-bootstrap}.
Scope: 30--40 pages.
\item \emph{Arithmetic Siegel--Weil at bad primes for \(g\geq 3\)},
extending Li--Liu 2021 arXiv:2101.09485. Scope: 40--60 pages.
\item \emph{Harder-Eisenstein-congruence prime at genus \(g\)} replacing
\(p=691\) in F6 (\Cref{v4-arith:conj:F6-chenevier-691}).
% RECTIFICATION-FLAG: Gate 1: The value p=2028 is not a prime. Harder's genus-2 congruences predict a prime p | Num(B_{10}/10) or similar; the tabulation of Bergstroem--Faber--van der Geer should be re-read to identify the correct prime congruence modulus for weight-(10,13) on Sp_4(Z).
Systematic tabulation by Bergstr\"om--Faber--van der Geer 2014
(\emph{Sel.~Math.}~20) identifies the relevant Hecke-eigenvalue congruences.
Scope: 20--30 pages.
\item \emph{Arithmetic \(\mathsf{B}\)-row genus-\(g\) ceiling.} F4
(\Cref{v4-arith:conj:F4-B-row}) sharpens from the single integer
\(\kappa^{\mathrm{Ar}}+(\kappa^{!})^{\mathrm{Ar}}=8\) to a
genus-dependent Shimura-signature-indexed sum, via Bruinier--Heegner
Chern-class reciprocity on \(\mathrm{Sh}(\mathrm{O}(2,2g-1))\) with
accidental isomorphism \(\mathrm{O}(2,3)\cong\mathrm{Sp}_4/\{\pm I\}\)
at \(g=2\). Scope: 30--50 pages.
\end{enumerate}

The total residual scope is approximately 375--605 pages, organised as
three to four paper-length treatments. Paper~(i): obligations 1--4 + 7
(Beilinson / Positselski / bar-cobar on 2-dim base plus Virasoro
bootstrap). Paper~(ii): obligations 5, 6 (\(GL_{2g}\) Fargues--Scholze
local plus global Arinkin--Gaitsgory). Paper~(iii): obligations 8, 10
(Kudla--Rapoport and Bruinier--Heegner at genus \(g\)). Paper~(iv):
obligation~9 and consolidation (Harder-Eisenstein congruences; all-genus
F-RH corollary).

\section{Orthogonal and Compatible Windows}
\label{v4-arith:ag:orthogonal}

The six windows are compatible but mathematically distinct: each
captures a different facet of the datum \(\mathcal{D}_g\), and their
pairwise compatibility is an independent mathematical fact in each case.
Pairwise compatibility:

\begin{itemize}
\item \textbf{Windows~I and V} share the arithmetic surface as base and
agree on the weight-\(1\) Hasse--Weil functional equation. They differ on
operadic dimension: Window~I keeps the archimedean \(E_1\)-chiral structure;
Window~V lifts to \(E_2\)-chiral via Dunn. Compatibility is the content of
Chapter~\ref{v4-arith:arakelov-bar}'s 2-dim extension.
\item \textbf{Windows~II and V} differ on arithmetic scheme dimension
(1-dim vs.\ 2-dim). Window~II encodes the genus as coefficient enhancement
of the base field (totally imaginary \(K\) of degree \(2g\)); Window~V keeps
\(K=\mathbb{Q}\) but lifts the base scheme to 2-dim via Dunn additivity.
Compatibility passes through Tate 1950 adelic Fourier at number-field level.
\item \textbf{Windows~III and VI} are spectral: Window~III gives the local
\((\varphi,\Gamma)\)-module moduli; Window~VI gives the global automorphic
side. Compatibility via Fargues--Scholze arXiv:2102.13459 local Langlands.
\item \textbf{Window~IV} is moduli-theoretic: \(\mathcal{T}_g^{\mathrm{arith}}\)
carries the Hodge--Arakelov decoration on the theta-generic locus, and
its point-by-point inputs supply Window~I. Compatibility via the Torelli
map: a genus-\(g\) curve's Jacobian is a point of
\(\mathcal{X}_{2g}^{\mathrm{crys,ord}}\) (Window~III) via its \(F\)-crystal,
and a point of the Siegel modular threefold (Window~VI) via its
polarisation.
\end{itemize}

Orthogonal to these six windows: Deninger's \(\mathbb{F}_1\) programme
(motivational context for the arithmetic-dynamical interpretation of
\(\mathrm{Spec}\,\mathbb{Z}\), recorded in Chapter~\ref{v4-arith:manifesto});
the Caraiani--Scholze unitary-Shimura programme at genera where no
Shimura model exists (moduli of curves \(\mathcal{M}_g\) is not Shimura for
\(g\geq 4\)); the Bethe-algebra heuristic presentation, which operates at the
zero-Fock per-place level and does not carry \(g\)-dependence directly.

The identity \(\kappa^{\mathrm{Ar},g}+(\kappa^{!})^{\mathrm{Ar},g}=0\)
of \Cref{v4-arith:ag:thm:all-genus} holds at every \(g\geq 0\). The
F-RH axis of Chapter~\ref{v4-arith:rh-reformulation} is the \(g=0\)
fibre of the rank-graded family whose other fibres are the Hasse--Weil
Riemann Hypothesis for Jacobians of curves of genus \(g\), the
generalised Riemann Hypothesis for Dedekind zeta of totally imaginary
fields of degree \(2g\), and the Siegel-spin Riemann Hypothesis for
cusp forms on \(\mathrm{GSp}_{2g}\). The residual obligations of
\Cref{v4-arith:ag:residual} quantify what must be proved to convert the
conditional biconditional into an unconditional one.

Chapter~\ref{v4-arith:3d-acs-E1bar} extends the arithmetic-modular axis
into arithmetic three-dimensional Chern--Simons theory. The conformal
vector implicit in \Cref{v4-arith:ag:thm:all-genus} sources an
arithmetic Kac--Moody enhancement of \(\mathcal{V}^{\mathrm{prim}}\)
(\Cref{v4-arith:3d-acs:thm:km-lift}), which in turn sources a
finite-level arithmetic Chern--Simons partition sum and a conjectural
renormalised \(GL_2(\widehat{\mathbb{Z}})\) limit on the \'etale
3-space \(\overline{\mathrm{Spec}\,\mathbb{Z}}\)
(\Cref{v4-arith:3d-acs:thm:finite-level-ACS},
\Cref{v4-arith:3d-acs:conj:qu-acs}). The equalities
\[
Z_{\mathrm{ACS}^{\mathrm{qu}}}=\xi(k^{\mathrm{Ar}}),
\qquad
\partial\mathrm{ACS}^{\mathrm{qu}}=\mathcal{V}^{\mathrm{prim}}
\]
are conditional on profinite renormalisation and the local trace
comparison (\Cref{v4-arith:3d-acs:thm:bulk-boundary-Vprim}). The
ordered incidence model is computed by the normalised ordered bar
construction on the auxiliary place order \(\mathbb{P}\cup\{\infty\}\)
(\Cref{v4-arith:3d-acs:thm:chi-hom-bar}); the comparison with
full Beilinson--Drinfeld / Costello--Gwilliam arithmetic chiral homology
is conditional (\Cref{v4-arith:3d-acs:cor:BDCG-ordered-comparison}).
F-RH becomes a conditional spectral theorem on the partition-function
zero locus (\Cref{v4-arith:3d-acs:thm:arith-gravity-FRH}); the
all-genus extension
(\Cref{v4-arith:3d-acs:cor:all-genus-arith-gravity}) covers every row
of the arithmetic-modular family.
